\documentclass[../../main.tex]{subfiles}
% \documentclass[a4paper,10pt]{ltjsarticle}
\begin{document}
\rhead{京都大学後期$1997$年$5$番}
\setcounter{equation}{0}
\begin{oframed}
  \subfile{problem}
\end{oframed}
\setlength{\columnseprule}{0.4pt}
\begin{multicols}{2}
  {\bf [解答]}

  \begin{figure}[H]
    \centering
    \begin{tikzpicture}[>=stealth,every node/.style={font=\small}]
      \node[draw=blue!85!black,rounded corners=2pt,align=center,
      minimum width=3.4cm,minimum height=1.0cm] (before) at (0,0)
      {$k$回目直前\\青$B_k$・赤$R_k$・黄$Y_k$};
      \node[draw=teal!85!black,rounded corners=2pt,align=center,
      minimum width=3.4cm,minimum height=1.0cm] (after) at (0,-2.3)
      {$k$回目終了\\合計$a+b+c+k$枚};
      \draw[->,very thick,blue!85!black] (before) -- (after)
      node[midway,right] {引いた色を1枚増やす};
      \node[draw=orange!90!black,rounded corners=2pt,align=center,
      minimum width=3.5cm] (added) at (0,-4.4)
      {$k+1$回目に直前の追加カードを引く};
      \node[draw=gray!80!black,rounded corners=2pt,align=center,
      minimum width=3.5cm] (old) at (0,-6.3)
      {$k+1$回目にそれ以外を引く};
      \draw[->,thick,orange!90!black] (after) -- (added)
      node[midway,right] {$\frac{1}{a+b+c+k}$};
      \draw[->,thick,gray!70!black] (after.west)
      to[out=200,in=160] node[left] {$\frac{a+b+c+k-1}{a+b+c+k}$} (old.west);
    \end{tikzpicture}
    \caption{カードの追加と$k+1$回目の分岐}
    \label{1997-5:fig:1}
  \end{figure}

  初期状態で青のカード$a$枚，赤のカード$b$枚，黄のカードが$c$枚であるとする．
  この時，$k$回目に青のカードをとり出す確率を$Q_k$とする．
  $Q_{k+1}$について，$k+1$回目にとり出すカードが$k$回目に加えたものか否かで場合分けする．
  $k$回目が終わった時点で，カードは全部で$a+b+c+k$枚存在する．
  $k+1$回目に引くカードが$k$回目に加えたものであるとき，それが青である確率は定義から$Q_k$である．
  一方$k+1$回目に引くカードが$k$回目に加えたものでないとき，それら$a+b+c+k-1$枚から青を引く確率は状況が$k$回目の時と全く同じため$Q_k$である．
  従って，漸化式は
  \begin{align}
    Q_{k+1}
    & = \dfrac{1}{a+b+c+k} Q_k + \dfrac{a+b+c+k-1}{a+b+c+k} Q_k \\
    & = Q_k
  \end{align}
  である．
  この漸化式をくり返し用いて，任意の自然数$n$に対して
  \begin{align}
    Q_n = Q_1 = \dfrac{a}{a+b+c}
  \end{align}
  を得る．従って，ここに$(a,b,c)=(3,2,1)$を代入して
  \begin{align}
    Q_n=\dfrac{1}{2} \label{eq:1}
  \end{align}
  を得る．

  以下，題意の期待値を求める．
  $n$回目の試行が完了した時の青の個数は，最初の$3$個と，途中で引いた青の数の和に等しい．
  まず，以下のように確率変数$X_k$を定める．
  \begin{align}
    X_k =
    \begin{dcases}
      1 & (k\text{回目に青を引く})  \\
      0 & (\text{otherwise})
    \end{dcases}
  \end{align}
  \cref{eq:1}より，$X_k$の期待値$E(X_k)$は
  \begin{align}
    E(X_k) = Q_k = \dfrac{1}{2}
  \end{align}
  である．
  求める期待値$E_n$は
  \begin{align}
    E_n = 3 + E(X_1+X_2+\cdots+X_n)
  \end{align}
  である．期待値の加法定理から
  \begin{align}
    E(n)
    & = 3+E(X_1+\dots+X_n)         \\
    & = 3+\sum_{k=1}^n E(X_k)      \\
    & = 3+\sum_{k=1}^n \dfrac{1}{2} \\
    & = \dfrac{n}{2}+3
  \end{align}
  が求める答えである．$\cdots$(答)
  {\bf [解説]}

  \newpage
\end{multicols}
\end{document}
